% `tex.web` §317: an error context is *bounded*, and this is the whole rule.
%
%     if l+first_count<=half_error_line then
%       begin p:=0; n:=l+first_count; end
%     else begin print("..."); p:=l+first_count-half_error_line+3;
%       n:=half_error_line; end;
%     for q:=p to first_count-1 do print_char(trick_buf[q mod error_line]);
%     print_ln;
%     for q:=1 to n do print_char(" ");
%     if m+n<=error_line then p:=first_count+m
%     else p:=first_count+(error_line-n-3);
%     for q:=first_count to p-1 do print_char(trick_buf[q mod error_line]);
%     if m+n>error_line then print("...")
%
% The head — `l.4 `, `<to be read again> `, a macro's name — is printed before
% the pseudoprint begins and is never itself cut. What gives is the *tail* of
% what has been read, marked with a leading `...`, and the *head* of what has
% not, marked with a trailing one. So the first line comes to at most
% `half_error_line` and the second to at most `error_line`.
%
% Neither is a TeX parameter; a document can neither read nor set them, so they
% are constants and this case is how their values were established rather than
% assumed: 50 and 79, the numbers TeX Live builds with.
%
% Every line here is long on purpose. Without the rule ours printed the source
% entire, which is a third of trip.log's error contexts differing in a way that
% reads as wrong content rather than as an unbounded line.
\catcode`\{=1 \catcode`\}=2 \catcode`\#=6

% Both halves at once: the read part overflows to the left, the unread part to
% the right.
\count1=\undefinedcontrolsequencewithaveryverylongname AAAAAAAAAA BBBBBBBBBB CCCCCCCCCC DDDDDDDDDD EEEEEEEEEE FFFFFFFFFF
\message{PROBE:one=\the\count1}

% §319 prints only a macro's *name* before the pseudoprint starts: its
% parameter text and the `->` belong to the list being shown, and so are part
% of what gets cut. Keeping them in the head breaks a truncated macro context
% two characters late.
\def\aaaaaaaaaaaaaaaaaaaaaaaaaaaaaaaaaaa#1{\count2=\undefinedtwo #1XXXXXXXXXXXXXXXXXXXXXXXXXXXXXXXXXXXXXXXXXXXXXXXXXXXXXXXXXXXXXX}
\aaaaaaaaaaaaaaaaaaaaaaaaaaaaaaaaaaa{Y}
\message{PROBE:two=\the\count2}

% A short context is left alone — the bound must not fire where TeX does not
% reach it.
\count3=\undefthree
\message{PROBE:three=\the\count3}
\end
